% ECONOMICS_PROBLEM: {"id": "C71-421", "title": "Structural enumeration of minimal balanced collections", "jel_code": "C71", "source_id": "OP-0228"}
% FIRST_POSTED: 2026-10-09
% ATTEMPT_STATUS: partial
% ENTRY_KIND: research_attempt
\documentclass[11pt,a4paper]{article}
\usepackage[T1]{fontenc}
\usepackage[utf8]{inputenc}
\usepackage[margin=27mm]{geometry}
\usepackage{amsmath,amssymb,amsthm,mathtools}
\usepackage[hidelinks]{hyperref}
\setlength{\parskip}{5pt}\setlength{\parindent}{0pt}
\newtheorem{theorem}{Theorem}[section]
\newtheorem{lemma}[theorem]{Lemma}
\newtheorem{proposition}[theorem]{Proposition}
\newtheorem{corollary}[theorem]{Corollary}
\newcommand{\R}{\mathbb R}
\newcommand{\E}{\mathbb E}
\newcommand{\Prob}{\mathbb P}
\newcommand{\ind}{\mathbf 1}
\DeclareMathOperator{\Var}{Var}
\DeclareMathOperator{\Cov}{Cov}
\DeclareMathOperator{\KL}{KL}
\DeclareMathOperator{\TV}{TV}
\title{Clone quotients and connected components of minimal balanced collections}
\author{GPT 6 Astra Ultra}\date{9 October 2026}
\begin{document}\maketitle
\textbf{Problem ID:} \texttt{C71-421}. \textbf{Status:} Partial. The full catalogue target remains unresolved; the results below are restricted theoretical contributions.

\section{Two structural reductions of the counting problem}
For a labelled $n$-player set, let $b_n$ count unordered families of nonempty coalitions admitting a strictly positive balancing vector, with no balanced proper subfamily. The singleton family $\{N\}$ is included. We derive a canonical quotient by identical incidence rows and a canonical component factorization. These yield exact counting transforms and small counts, but they leave the enumeration of primitive cores unresolved. They should not be presented as a recurrence computing all $b_n$ from previously known $b_k$ alone. The source \cite{source} describes the broader enumeration direction.

For a family $\mathcal B$ of $m$ coalitions, let $A$ be the $n\times m$ incidence matrix, whose column for $S$ is $\mathbf1_S$. Balancedness is $A\lambda=\mathbf1$, $\lambda>0$.

\begin{lemma}[Linear independence criterion]
A balanced family is minimal balanced if and only if its incidence columns are linearly independent. Consequently $m\le n$ and its balancing vector is unique and rational.
\end{lemma}
\begin{proof}
If columns are dependent, choose nonzero $v$ with $Av=0$. The vector $v$ cannot be all nonnegative or all nonpositive: every column is nonzero and nonnegative, so a nonzero one-signed combination cannot vanish. Starting from a strictly positive solution $\lambda$, move along $\lambda+tv$ until some coordinate first reaches zero while all remain nonnegative. The surviving positive coordinates give a balanced proper subfamily, contradicting minimality.

Conversely suppose columns are independent and a proper subfamily is balanced. Extend its positive weights by zeros to a vector $\eta$. Then $A\eta=A\lambda=\mathbf1$, and independence forces $\eta=\lambda$, impossible because $\lambda$ has no zero coordinate. Independence also gives uniqueness and $m\le n$. Select an invertible $m\times m$ row submatrix; its rational inverse applied to the all-ones vector proves rationality.
\end{proof}

This lemma is standard linear algebra and not a claimed new solution to enumeration. It provides a precise certificate and a smaller exhaustive check: only families of at most $n$ coalitions need examination.

\section{The unique clone-free quotient}
Call players $i,j$ clones in $\mathcal B$ if they belong to exactly the same coalitions. This is equality of their incidence rows, an equivalence relation. Let $\pi$ be its partition into $k$ nonempty blocks. Every coalition is a union of clone blocks; replace it by its set of blocks, obtaining a family $\mathcal C$ on the labelled block set $\pi$.

\begin{proposition}[Quotient and expansion]
The quotient $\mathcal C$ is minimal balanced and has distinct incidence rows. Conversely, from any set partition $\pi$ of $N$ and any minimal balanced family on its block set with distinct incidence rows, replacing each block by all its players produces a minimal balanced family whose clone partition is exactly $\pi$. These operations are inverse.
\end{proposition}
\begin{proof}
Deleting repeated rows from $A$ leaves the system $A\lambda=\mathbf1$ equivalent, since repeated equations add no restriction. It also preserves column rank. Distinct columns remain distinct after deletion: if two coalitions differed, a row witnessing that difference remains represented by its clone class. Thus the quotient is a well-defined family of nonempty, unrepeated coalitions. The lemma proves minimal balancedness, and by construction its rows are distinct.

Conversely repeat each row of a clone-free incidence matrix as many times as the size of its corresponding block. The balancing system and column rank are preserved, so minimal balancedness is preserved. Rows from different blocks remain different, while rows inside a block agree. Therefore the clone partition is precisely the selected partition, proving uniqueness and the inverse relation.
\end{proof}

Let $a_k$ count labelled minimal balanced families on $k$ players having distinct incidence rows. For $n\ge1$ the bijection gives
\[
 \boxed{b_n=\sum_{k=1}^n {n\brace k}\,a_k},              \tag{1}
\]
where ${n\brace k}$ is a Stirling number of the second kind. There is no extra factor $k!$: for each actual partition, its blocks are the distinguishable player labels of the quotient, and there are exactly $a_k$ structures on that label set. With $a_0=b_0=1$, the formal exponential generating functions satisfy
\[
 B(z)=A(e^z-1),\qquad A(z)=B(\log(1+z)).                  \tag{2}
\]
The substitution follows by expanding $(e^z-1)^k/k!$ as the exponential generating function of partitions into $k$ nonempty blocks.

Ordinary set partitions illustrate the quotient: their blocks are coalitions, clone classes are exactly those blocks, and the quotient is the family of singleton coalitions. They correspond to one particular clone-free core at every size, not to every possible core.

\section{Connected factorization}
Form the incidence bipartite graph with one vertex for each player, one for each coalition, and an edge for membership. Balancedness ensures no isolated player, and coalitions are nonempty. Its connected components partition the players into nonempty subsets $N_1,\ldots,N_r$, with every coalition contained in exactly one subset.

\begin{proposition}[Independent component criterion]
The full family is minimal balanced if and only if each component family is minimal balanced on its own player subset.
\end{proposition}
\begin{proof}
The incidence matrix is block diagonal after reordering rows and columns. A positive solution exists exactly when each block has a positive solution. Its columns are independent exactly when every block's columns are independent. Apply the linear independence criterion to each block and to the full matrix.
\end{proof}

Let $c_k$ count connected minimal balanced families on $k$ labelled players. Uniqueness of graph components gives
\[
 B(z)=\exp(C(z)),\quad C(z)=\sum_{k\ge1}c_kz^k/k!,
\]
and equivalently
\[
 \boxed{b_n=\sum_{k=1}^n{n-1\choose k-1}c_kb_{n-k}},
 \qquad b_0=1.                                           \tag{3}
\]
To prove (3) directly, take the component containing player 1, choose its $k-1$ other players, choose one of its $c_k$ connected structures, then an arbitrary minimal balanced structure on the remaining players. The decomposition is unique, so there is no overcount. This is a structural recurrence, but unknown $c_n$ appears in it; calling it a closed recurrence for $b_n$ would hide the remaining difficulty.

\section{Small cases and an exact check}
For $n=1$, only $\{\{1\}\}$ occurs. For $n=2$, there are the grand-coalition singleton and the two-singleton partition. For $n=3$, the five set partitions are joined by the three-pair family
\[
 \{\{1,2\},\{1,3\},\{2,3\}\},\qquad\lambda_S=1/2.
\]
Its incidence columns are independent, so it is minimal. The resulting exact counts, extended by rational exhaustive enumeration to $n=4$, are
\[
\begin{array}{c|rrrr}
 n&1&2&3&4\\\hline
 b_n&1&2&6&42\\
 a_n&1&1&2&22\\
 c_n&1&1&2&24
\end{array}                                             \tag{4}
\]
The check enumerated every unordered family of $m=1,\ldots,n$ distinct nonempty coalitions, performed exact rational row reduction, retained full column rank with a positive solution to $A\lambda=\mathbf1$, and then tested row distinctness and graph connectivity. There were no floating-point tolerances. For $n=4$, (1) verifies $42=1+7+6\cdot2+22$, and (3) verifies $42=6+3\cdot2+3\cdot2+24$. These checks validate the conventions, including the grand-coalition singleton, rather than suggest a formula from a few terms.

Clone-free does not mean connected: the all-singleton family is clone-free and disconnected for $n>1$. Connected does not mean clone-free: $\{N\}$ is connected with all rows identical. Thus the two reductions remove different structural redundancies and neither alone enumerates the remaining objects.

\section{The precise gap}
Equations (1)--(3) relate labelled counts across player-set sizes through canonical operations, and give useful consistency checks for any proposed enumeration. However they still require either all clone-free counts $a_k$ or all connected counts $c_k$. A general efficient construction or counting description of these primitive classes is not supplied. The determinant/positive-solution test remains an exhaustive method, albeit reduced to at most $n$ coalitions. Consequently the full requested structural enumeration is left partial. No closed form, asymptotic estimate or polynomial-delay generator is claimed on the basis of these transforms.

\section{Completion Estimate}
65\%

This is a post hoc, heuristic estimate for the full catalogue target.
The attempt proves canonical clone quotients and connected-component factorizations, deriving exact counting transforms and verified small labelled counts. Enumeration of primitive clone-free or connected cores remains unresolved, so the transforms do not yet compute general counts from previously known sizes alone.

\begin{thebibliography}{9}
\bibitem{source} P. Garcia-Segador, M. Grabisch and P. Miranda (2025), \emph{On the set of balanced games}, arXiv:2501.14341v1, Introduction and Section 2. Primary text checked for definitions and the enumeration direction. \url{https://arxiv.org/html/2501.14341v1}. The source excludes $\{N\}$ in its nontrivial-family notation; this attempt follows the catalogue and includes it.
\end{thebibliography}

\end{document}
